% Bagian: computability
% Bab: recursive-functions
% Subbagian: pr-functions

\documentclass[../../../include/open-logic-section]{subfiles}

\begin{document}

\olfileid[id]{cmp}{rec}{prf}
\olsection{Fungsi Rekursif Primitif}

Mari kita catat kembali bagaimana kita dapat mendefinisikan fungsi baru dari
fungsi yang sudah ada dengan menggunakan rekursi primitif dan komposisi.

\begin{defn}
  \ollabel{defn:primitive-recursion} Andaikan $f$ merupakan fungsi berargumen
  $k$ ($k\ge 1$) dan $g$ merupakan fungsi berargumen $(k+2)$. Fungsi yang
  \emph{didefinisikan dengan rekursi primitif dari $f$ dan $g$} adalah fungsi
  berargumen $(k+1)$, yaitu~$h$, yang didefinisikan oleh persamaan-persamaan
  \begin{align*}
    h(x_0,\dots,x_{k-1},0) & =  f(x_0,\dots,x_{k-1}) \\
    h(x_0,\dots,x_{k-1},y+1) & = g(x_0,\dots,x_{k-1}, y, h(x_0,\dots,x_{k-1}, y))
  \end{align*}
\end{defn}

\begin{defn}
  \ollabel{defn:composition} Andaikan $f$ merupakan fungsi berargumen $k$,
  dan $g_0$, \dots, $g_{k-1}$ merupakan $k$ fungsi yang semuanya berargumen
  $n$. Fungsi yang \emph{didefinisikan dengan komposisi dari $f$ dan $g_0$,
    \dots,~$g_{k-1}$} adalah fungsi berargumen $n$, yaitu~$h$, yang
  didefinisikan oleh
  \[
  h(x_0, \dots, x_{n-1}) =
  f(g_0(x_0, \dots, x_{n-1}), \dots, g_{k-1}(x_0, \dots, x_{n-1})).
  \]
\end{defn}

Selain $\Succ$ dan fungsi proyeksi
\[
\Proj{n}{i}(x_0,\dots,x_{n-1}) = x_i,
\]
untuk setiap bilangan asli $n$ dan $i < n$, kita akan memasukkan fungsi
$\Zero(x) = 0$ ke dalam golongan fungsi rekursif primitif.

\begin{defn}
  Himpunan fungsi rekursif primitif adalah himpunan fungsi dari $\Nat^n$ ke
  $\Nat$ yang didefinisikan secara induktif oleh klausa-klausa berikut:
  \begin{enumerate}
  \item $\Zero$ bersifat rekursif primitif.
  \item $\Succ$ bersifat rekursif primitif.
  \item Setiap fungsi proyeksi $\Proj{n}{i}$ bersifat rekursif primitif.
  \item Jika $f$ merupakan fungsi rekursif primitif berargumen $k$ dan $g_0$,
    \dots,~$g_{k-1}$ merupakan fungsi rekursif primitif berargumen $n$, maka
    komposisi $f$ dengan $g_0$, \dots,~$g_{k-1}$ bersifat rekursif primitif.
  \item Jika $f$ merupakan fungsi rekursif primitif berargumen $k$ dan $g$
    merupakan fungsi rekursif primitif berargumen $k+2$, maka fungsi yang
    didefinisikan dengan rekursi primitif dari $f$ dan $g$ bersifat rekursif
    primitif.
\end{enumerate}
\end{defn}

\begin{explain}
Secara lebih ringkas, himpunan fungsi rekursif primitif adalah himpunan
terkecil yang memuat $\Zero$, $\Succ$, dan fungsi proyeksi~$\Proj{n}{j}$,
serta tertutup terhadap komposisi dan rekursi primitif.

Cara lain untuk menggambarkan himpunan fungsi rekursif primitif adalah dengan
mendefinisikannya berdasarkan ``tahap.'' Misalkan $S_0$ menyatakan himpunan
fungsi awal: $\Zero$, $\Succ$, dan proyeksi. Fungsi-fungsi ini adalah fungsi
rekursif primitif tahap~$0$. Setelah tahap $S_i$ didefinisikan, misalkan
$S_{i+1}$ adalah tahap sebelumnya bersama himpunan seluruh fungsi yang diperoleh dengan
menerapkan satu komposisi atau satu rekursi primitif pada fungsi-fungsi yang
sudah berada dalam~$S_i$. Maka
\[
S = \bigcup_{i \in \Nat} S_i
\]
adalah himpunan semua fungsi rekursif primitif.
\end{explain}

Mari kita pastikan bahwa $\Add$ merupakan fungsi rekursif primitif.

\begin{prop}
  Fungsi penjumlahan $\Add(x,y) = x+y$ bersifat rekursif primitif.
\end{prop}

\begin{proof}
Kita sudah mempunyai definisi rekursif primitif bagi $\Add$ dalam bentuk dua
fungsi $f$ dan~$g$ yang sesuai dengan pola
\olref{defn:primitive-recursion}:
\begin{align*}
  \Add(x_0, 0) & = f(x_0) = x_0\\
  \Add(x_0, y+1) & = g(x_0, y, \Add(x_0, y)) =
  \Succ(\Add(x_0, y))
\end{align*}
Jadi, $\Add$ bersifat rekursif primitif asalkan $f$ dan $g$ juga demikian.
$f(x_0) = x_0 = \Proj{1}{0}(x_0)$, dan fungsi proyeksi termasuk rekursif
primitif, sehingga $f$ bersifat rekursif primitif. Fungsi $g$ adalah fungsi
berargumen tiga, yaitu $g(x_0, y, z)$, yang didefinisikan oleh
\[
g(x_0, y, z) = \Succ(z).
\]
Hal ini belum menunjukkan bahwa $g$ bersifat rekursif primitif karena $g$ dan
$\Succ$ bukan fungsi yang persis sama: $\Succ$ berargumen satu, sedangkan $g$
harus berargumen tiga. Namun, kita dapat mendefinisikan $g$ secara ``resmi''
dengan komposisi sebagai
\[
g(x_0, y, z) = \Succ(\Proj{3}{2}(x_0, y, z)).
\]
Karena $\Succ$ dan $\Proj{3}{2}$ termasuk fungsi rekursif primitif, $g$ juga
bersifat rekursif primitif, sebab fungsi itu dapat didefinisikan dengan
komposisi dari fungsi-fungsi rekursif primitif.
\end{proof}

\begin{prop}
  \ollabel{prop:mult-pr}
  Fungsi perkalian $\Mult(x,y) = x \cdot y$ bersifat rekursif primitif.
\end{prop}

\begin{proof}
  Latihan.
\end{proof}

\begin{prob}
  Buktikan \olref[cmp][rec][prf]{prop:mult-pr} dengan menunjukkan bahwa
  definisi rekursif primitif bagi $\Mult$ dapat ditulis dalam bentuk yang
  disyaratkan oleh \olref[cmp][rec][prf]{defn:primitive-recursion}, dan bahwa
  fungsi $f$ serta~$g$ yang bersesuaian bersifat rekursif primitif.
\end{prob}

\begin{ex}
Berikut contoh pertama kita tentang definisi rekursif primitif:
\begin{align*}
h(0) & =  1 \\
h(y+1) & =  2 \cdot h(y).
\end{align*}
Fungsi ini tidak dapat dimasukkan ke dalam bentuk yang disyaratkan oleh
\olref{defn:primitive-recursion} karena $k=0$. Definisi tersebut juga memuat
konstanta $1$ dan $2$. Untuk mengatasi masalah pertama, mari kita tambahkan
argumen semu dan definisikan fungsi~$h'$:
\begin{align*}
h'(x_0, 0) & =  f(x_0) = 1 \\
h'(x_0, y+1) & =  g(x_0, y, h'(x_0, y)) = 2 \cdot h'(x_0, y).
\end{align*}
Fungsi $f(x_0) = 1$ dapat didefinisikan dari $\Succ$ dan $\Zero$ dengan
komposisi: $f(x_0) = \Succ(\Zero(x_0))$. Fungsi $g$ dapat didefinisikan dengan
komposisi dari $g'(z) = 2 \cdot z$ dan proyeksi:
\begin{align*}
g(x_0, y, z) & = g'(\Proj{3}{2}(x_0, y, z))
\intertext{dan selanjutnya $g'$ dapat didefinisikan dengan komposisi sebagai}
g'(z) & = \Mult(g''(z), \Proj{1}{0}(z))
\intertext{serta}
g''(z) & = \Succ(f(z)),
\end{align*}
dengan $f$ seperti di atas: $f(z) = \Succ(\Zero(z))$. Setelah memperoleh~$h'$,
kita dapat kembali menggunakan komposisi untuk menetapkan $h(y) =
h'(\Proj{1}{0}(y),\Proj{1}{0}(y))$. Hal ini menunjukkan bahwa $h$ dapat
didefinisikan dari fungsi-fungsi dasar dengan serangkaian komposisi dan rekursi
primitif, sehingga $h$ bersifat rekursif primitif.
\end{ex}

\end{document}
